\section{模曲线与模空间}
\label{sec:modular-curves-moduli}
% ============================================================================

\textbf{为什么需要模曲线？}

\cref{cor:torus-tau} 告诉我们：$Y(1) = \SL_2(\Z) \backslash \HH$
参数化所有椭圆曲线的同构类，通过 $j$-不变量与 $\C$ 同构。
这很好——但有时我们不只想分类椭圆曲线，
还想分类"带有额外数据"的椭圆曲线。

\textbf{典型的"额外数据"是什么？}
以 Wiles 证明 FLT 为例。核心步骤是：把"级别 $N$ 的模形式"与"有理椭圆曲线"对应起来。
这里"级别 $N$"就意味着椭圆曲线上记录了某种 $N$ 阶结构（具体是一个 $N$ 阶循环子群）。
要参数化"带 $N$ 阶子群的椭圆曲线"，$Y(1)$ 就不够用了——它把有无额外结构的曲线混为一谈。

\textbf{解决方案：用更小的对称群。}
如果我们只用 $\Gamma_0(N)$ 的等价关系来识别椭圆曲线
（即只把"$N$ 阶子群相同"的曲线视为等价），
就得到一个更细的分类空间 $Y_0(N) = \Gamma_0(N) \backslash \HH$。
不同的 $N$ 和不同的子群类型（$\Gamma_0, \Gamma_1, \Gamma$）
给出不同"粗细程度"的分类。

这些分类空间——\textbf{模曲线}——不只是集合，
它们还具有代数曲线的结构（将在\cref{ch:riemann-surfaces}中看到这是紧黎曼面），
可以定义在有理数甚至整数上，成为研究椭圆曲线算术的核心工具。

\textbf{谷山-志村猜想}（现在是定理）断言：每条有理椭圆曲线 $E$ 都是某个模曲线 $X_0(N)$ 的"商"，
即存在有理映射 $X_0(N) \to E$。这就是 FLT 证明的骨架。

\begin{definition}[模曲线]
\label{def:modular-curves}
对同余子群 $\Gamma$，定义\textbf{（开）模曲线}
\[
  Y(\Gamma) = \Gamma \backslash \HH.
\]
特别地，$Y_0(N) = \GammaO(N) \backslash \HH$，
$Y_1(N) = \GammaI(N) \backslash \HH$，
$Y(N) = \Gamma(N) \backslash \HH$。

通过添加尖点得到\textbf{（紧化）模曲线}
\[
  X(\Gamma) = \Gamma \backslash \HH^*.
\]
$X(\Gamma)$ 是紧黎曼面（将在\cref{ch:riemann-surfaces}中详细讨论）。
\end{definition}

\begin{theorem}[模诠释]
\label{thm:moduli-interpretation}
模曲线参数化带附加结构的椭圆曲线（在 $\C$ 上）：
\begin{enumerate}
  \item $Y(1) = \SL_2(\Z) \backslash \HH$ 参数化复环面的同构类
    （\cref{cor:torus-tau}）。通过 $j$-不变量，$Y(1) \cong \C$。
  \item $Y_0(N)$ 参数化 $\{(E, C)\}$ 的同构类，其中 $E$ 是椭圆曲线，
    $C \subset E$ 是 $N$ 阶循环子群（等价于核为 $\Z/N\Z$ 的同源
    $E \to E/C$）。
  \item $Y_1(N)$ 参数化 $\{(E, P)\}$ 的同构类，其中 $P \in E$ 是
    $N$ 阶点。
  \item $Y(N)$ 参数化 $\{(E, (P, Q))\}$ 的同构类，其中 $(P, Q)$ 是
    $E[N]$ 的一组有序基。
\end{enumerate}
\end{theorem}

\begin{proof}[解释]
\textbf{(2)}：设 $\tau \in \HH$，对应环面 $E_\tau = \C/\Lambda_\tau$
（$\Lambda_\tau = \Z\tau + \Z$）。取 $C = \langle 1/N \rangle \subset E_\tau$，
即 $C = \frac{1}{N}\Z/\Z$ 是 $N$ 阶循环子群。

何时 $(E_\tau, C_\tau) \cong (E_{\tau'}, C_{\tau'})$？
需要 $\alpha \in \C^*$ 使得 $\alpha\Lambda_\tau = \Lambda_{\tau'}$
且 $\alpha C_\tau = C_{\tau'}$。写出条件：
$\alpha\tau = a\tau' + b$，$\alpha = c\tau' + d$（$ad - bc = 1$），
且 $\alpha \cdot \frac{1}{N} \equiv \frac{j}{N} \pmod{\Lambda_{\tau'}}$
对某 $j$ 与 $N$ 互素。

这等价于 $\frac{c\tau' + d}{N} \in \frac{1}{N}\Z + \Lambda_{\tau'}$，
即 $c \equiv 0 \pmod{N}$。
故等价关系正是 $\GammaO(N)$ 的作用。

\textbf{(3) 和 (4)}：类似分析。
\end{proof}

\subsection*{Hauptmodul 与亏格 $0$ 的函数域}

CSS 中反复强调：模曲线除了是参数空间，也有自己的函数域。对亏格 $0$ 的曲线，这个函数域最简单——它总是由一个函数生成。

\begin{definition}[Hauptmodul]
\label{def:hauptmodul}
\index{Hauptmodul}
设 $X$ 是定义在 $\C$ 上的紧黎曼面，且 $\genus(X)=0$。若一个亚纯函数
\[
  t\colon X \longrightarrow \mathbb{P}^1(\C)
\]
给出双有理同构，等价地说 $\C(X)=\C(t)$，则称 $t$ 为 $X$ 的\textbf{Hauptmodul}（主模）。
\end{definition}

\begin{proposition}[亏格 $0$ 曲线的函数域]
\label{prop:genus0-hauptmodul}
若 $X$ 是亏格 $0$ 的紧黎曼面，则 $X\cong \mathbb{P}^1(\C)$，从而存在 Hauptmodul $t$ 使得
\[
  \C(X)=\C(t).
\]

若 $P\in X$ 是一点，则可取 $t$ 在 $P$ 处有一个单极点、在别处全纯；这样的 $t$ 已经生成整个函数域。
\end{proposition}

\begin{proof}[证明思路]
由黎曼面分类，亏格 $0$ 的紧黎曼面与 $\mathbb{P}^1(\C)$ 双全纯同构。把该同构复合到坐标函数 $z$ 上，就得到一个亚纯函数 $t$。而 $\mathbb{P}^1$ 的亚纯函数正是有理函数，所以 $\C(X)$ 同构于 $\C(z)=\C(t)$。
\end{proof}

对模曲线来说，这个定义尤其重要：当 $X(\Gamma)$ 的亏格为 $0$ 时，我们就能用一个单参数来描述整条曲线，所有模函数都是这个参数的有理函数。

\begin{example}[$j$ 与若干低级别 Hauptmodul]
\label{ex:hauptmodul-examples}
\leavevmode
\begin{itemize}
  \item 对 $X(1)=\SL_2(\Z)\backslash \HH^*$，经典 $j$-不变量在尖点 $\infty$ 处有简单极点，且没有别的极点，因此
    \[
      \C(X(1))=\C(j).
    \]
    也就是说，$j(\tau)$ 正是 $X(1)$ 的 Hauptmodul。
  \item 对经典模曲线 $X(2)$，Hauptmodul 是 Legendre 的 $\lambda(\tau)$。它把带一组有序 $2$-扭点基的椭圆曲线写成
    \[
      y^2=x(x-1)(x-\lambda).
    \]
  \item 对 $X_0(2)$，一个常用的 Hauptmodul 是
    \[
      t_2(\tau)=\left(\frac{\eta(\tau)}{\eta(2\tau)}\right)^{24}=q^{-1}+O(1),
    \]
    其中 $\eta(\tau)=q^{1/24}\prod_{n\ge 1}(1-q^n)$ 是 Dedekind $\eta$ 函数。它满足
    \[
      j(\tau)=\frac{(t_2+256)^3}{t_2^2},
      \qquad
      j(2\tau)=\frac{(t_2+16)^3}{t_2}.
    \]
    因而 $j(\tau)$ 与 $j(2\tau)$ 都是 $t_2$ 的有理函数，说明 $\C(X_0(2))=\C(t_2)$。
  \item 对 $X_0(3)$，类似地可取
    \[
      t_3(\tau)=\left(\frac{\eta(\tau)}{\eta(3\tau)}\right)^{12}=q^{-1}+O(1),
    \]
    并有
    \[
      j(\tau)=\frac{(t_3+27)(t_3+243)^3}{t_3^3}.
    \]
    所以当 $X_0(N)$ 亏格为 $0$ 时，确实也存在一个单参数 Hauptmodul 来生成其函数域。
\end{itemize}
\end{example}

\subsection*{模点、同源类与扭曲}

前面已经从复环面角度说明了 $\SL_2(\Z)\backslash \HH$ 的意义。把这个结论改写成模曲线语言，就得到最基本的模诠释。

\begin{corollary}[复椭圆曲线的同构类]
\label{cor:complex-ec-moduli}
椭圆曲线 $E/\C$ 的同构类与 $Y(1)=\SL_2(\Z)\backslash \HH$ 的点一一对应；加入尖点后，$X(1)$ 只是把退化的结点三次曲线也补进去。

特别地，每个 $E/\C$ 都同构于某个 $E_\tau=\C/(\Z\tau+\Z)$，而且
\[
  E_\tau \cong E_{\tau'}
  \quad\Longleftrightarrow\quad
  \tau' = \gamma(\tau) \text{ 对某个 } \gamma\in \SL_2(\Z).
\]
\end{corollary}

\begin{proof}
这只是把 \cref{cor:torus-tau,thm:uniformization} 与模曲线记号合在一起重述。
\end{proof}

\begin{definition}[有理同源类]
\label{def:q-isogeny-class}
设 $E,E'/\Q$ 为椭圆曲线。若存在定义在 $\Q$ 上的同源 $E\to E'$，则称它们属于同一个\textbf{有理同源类}（isogeny class）。
\end{definition}

在模性定理里，真正与一个归一化新形式对应的通常不是单独一条椭圆曲线，而是一个 $\Q$-同源类：同一同源类中的曲线有相同的导子与同一个 Hasse--Weil $L$-函数，只是最小模型、挠点结构或 Tamagawa 数可能不同。

\begin{example}[两个特殊的 $j$ 值]
\label{ex:j0-j1728}
\leavevmode
\begin{itemize}
  \item $j=1728$：过 $\C$ 同构后，所有这样的椭圆曲线都可写成
    \[
      y^2=x^3-x,
    \]
    它有额外自同构 $z\mapsto iz$，所以
    \[
      \End(E_{\overline{\Q}})\supset \Z[i].
    \]
    过 $\Q$ 时，可以写成 $E_A\colon y^2=x^3+Ax$（$A\in \Q^\times$）；并且 $E_A\cong_{\Q}E_{A'}$ 当且仅当 $A'/A\in (\Q^\times)^4$。
  \item $j=0$：过 $\C$ 同构后，所有这样的椭圆曲线都可写成
    \[
      y^2=x^3-1,
    \]
    它有额外自同构 $z\mapsto \omega z$（$\omega=e^{2\pi i/3}$），所以
    \[
      \End(E_{\overline{\Q}})\supset \Z[\omega].
    \]
    过 $\Q$ 时，可以写成 $E_B\colon y^2=x^3+B$（$B\in \Q^\times$）；并且 $E_B\cong_{\Q}E_{B'}$ 当且仅当 $B'/B\in (\Q^\times)^6$。
\end{itemize}

这两个例子说明：同一个复同构类在 $\Q$ 上往往会裂成许多不同的有理同构类，而这些差别正是由扭曲控制的。
\end{example}

\begin{definition}[二次扭曲]
\label{def:quadratic-twist}
设
\[
  E\colon y^2=x^3+Ax+B
  \qquad (A,B\in \Q,\,4A^3+27B^2\neq 0)
\]
且 $d\in \Q^\times$。$E$ 关于 $d$ 的\textbf{二次扭曲}定义为
\[
  E^{(d)}\colon y^2=x^3+d^2Ax+d^3B.
\]
在 $\overline{\Q}$ 上，取 $\sqrt d$ 后映射
\[
  (x,y)\longmapsto (dx,d^{3/2}y)
\]
给出 $E\cong E^{(d)}$；但若 $d\notin (\Q^\times)^2$，则它们通常并不同构于 $\Q$。
\end{definition}

\textbf{注。} 扭曲的统一语言来自 Galois 上同调。固定 $E/\Q$ 后，所有在 $\overline{\Q}$ 上与 $E$ 同构、但在 $\Q$ 上可能不同构的模型，由集合
\[
  H^1\bigl(\Gal(\overline{\Q}/\Q),\Aut(E_{\overline{\Q}})\bigr)
\]
分类。若 $j(E)\neq 0,1728$，则 $\Aut(E_{\overline{\Q}})=\{\pm 1\}$，于是二次扭曲正对应于
\[
  H^1\bigl(\Gal(\overline{\Q}/\Q),\{\pm 1\}\bigr)
  \cong \Q^\times/(\Q^\times)^2.
\]
这就是为什么“按平方类扭曲”会自然出现。


% ============================================================================
